\chapter{Development Cycle and Trial Participants}
\label{Appendix6}

This appendix records who took part in the design and testing of uMatter
outside the two-person development team, and the weekly cycle through which
their input reached the product. The baseline platform evaluation that opened
the study is reported separately in Appendix~\ref{Appendix4}; the
Acknowledgments thank the same people by name.

\section{Weekly Review Cycle}
\label{sec:weekly-cycle}

Development ran in short phases, each delivering one coherent slice of the
platform (for example, the tracking and journaling features, the grant model
and the AI companion, or therapist booking and consultation). The requirements
in Section~\ref{sec:requirements} were not settled once at the start and then
implemented; they were revised phase by phase against three groups.

\begin{itemize}
  \item \textbf{The advisor.} Weekly supervision meetings with
    \tenGVHD{} covered scope, architecture, and what was realistic to
    deliver, and were where a proposed feature was first argued for or
    dropped.

  \item \textbf{Mental-health professionals.} The three professionals
    introduced in Chapter~\ref{Chapter1} reviewed both the features proposed
    for an upcoming phase and the features just completed, judging whether
    each was clinically sound and appropriate for adolescents.

  \item \textbf{Trial users.} The seven participants listed in
    Table~\ref{tab:participants} gave the first round of user-story interviews
    before implementation began, and were interviewed again against the
    working build as the phases were completed.
\end{itemize}

In any given week the team met either the professionals or the trial users,
depending on which set of proposed functions had been finished: a phase was
treated as closed only once its features had been shown to one of the two
groups, and the comments that came back were carried into the following phase.
The purpose of the cadence was to keep the delivered feature set close to what
the intended users and their clinicians actually asked for, rather than to what
the team assumed on their behalf. Several of the outcomes are visible in the
system: the periodic well-being surveys planned early on were dropped in favor
of the daily journal and 3-hourly mood check-ins after trial users called them
redundant, and the Treasure Box (Section~\ref{sec:most-valued}) exists because
a counsellor proposed reconnecting users with their reasons for living and the
CEO of LUMOS turned that idea into a concrete feature.

\section{User-Story and Feedback Participants}

Table~\ref{tab:participants} lists the trial users. All seven gave the first
round of user-story interviews and all seven returned for the closed testing
and feedback round, so their comments span the whole project rather than a
single snapshot of it. Two of them, Bảo Châu and Gia Phát, had also served as the two
evaluators of the baseline platform study in Appendix~\ref{Appendix4}.

\begin{table}[htbp]
  \centering
  \small
  \caption{Trial users who took part in the user-story and feedback rounds}
  \label{tab:participants}
  \begin{tabular}{|L{3.7cm}|L{6.3cm}|L{3.4cm}|}
    \hline
    \textbf{Participant} & \textbf{Status} & \textbf{Rounds} \\
    \hline
    Nguyễn Lê Bảo Châu & Psychology graduate, Ho Chi Minh City University of
      Education & User stories; feedback; baseline study \\
    \hline
    Nguyễn Phi Hùng & Business graduate & User stories; feedback \\
    \hline
    Lê Gia Phát & Business student, RMIT University Vietnam & User stories;
      feedback; baseline study \\
    \hline
    Trần Nhật Thanh & Student, Advanced Program in Computer Science,
      University of Science, VNU-HCM & User stories; feedback \\
    \hline
    Huỳnh Hữu Hậu & Student, Advanced Program in Computer Science,
      University of Science, VNU-HCM & User stories; feedback \\
    \hline
    Nguyễn Phước Ngọc Hương & Computer Science student, Ho Chi Minh City
      University of Technology, VNU-HCM & User stories; feedback \\
    \hline
    Trần Vĩ Quang & Computer Science student, Ho Chi Minh City University of
      Technology, VNU-HCM & User stories; feedback \\
    \hline
  \end{tabular}
\end{table}

The group is small and self-selected from the team's own circle, and it is
weighted toward university students rather than the secondary-school half of
the target population. It supports the qualitative findings reported in
Section~\ref{sec:expert-findings} and nothing stronger; no statistical claim is
made from it.

\section{Closed-Testing Cohort}

Publishing the Android client through Google Play required a closed test run
with a minimum number of opted-in testers over a continuous fourteen-day
period before the application could be promoted to production. Seventeen
testers took part in that track, recruited externally for the fourteen-day
window.

The role of this cohort was different from that of the trial users above.
Its members installed the published build and reported the defects they ran
into; they were not interviewed about the design and did not contribute to the
feature set. Their reports are the source of the fixes made during the release
phase, and are counted as testing effort rather than as user research.
